\subsection{Analýza nanopórových sekvenačných dát pomocou Dorado a Clover}
Najprv som stiahol potrebný dataset z verejného repozitára %Todo: doplniť odkaz na dataset
Potom som nan zavolal Dorado basecaller, ktorý som použil na prekladanie surových pod5 súborov do formátu BAM.
To som spravil pomocou nasledujúceho príkazu v termináli:
\begin{verbatim}        
dorado basecaller --input-format pod5 --output-format bam --input /cesta/k/datasetu --output /cesta/k/vystupu
\end{verbatim}
tam som si musel stiahnuť a nainštalovat Dorado, potom na zaklade pod5 datasetu som stiahnut dorado triedu, ktorá je určená na spracovanie daného datasetu. 
Na zaklade toho cez aky Nanopore bol dataset získaný.
po spracovaní Dorado basecallera som získal BAM súbor, ktorý obsahoval preložené sekvencie po 1000000 sekvencií dlkzy 1000000
To trvalo približne X hodín, v závislosti od veľkosti datasetu a výpočtovej kapacity použitého hardvéru. %
Následne som použil algoritmus Clover, ktorý zoskupil sekvencie do jednotlivých klastrov na základe ich podobnosti.
Problem je ze Clover nebere ako vstup BAM súbory, ale potrebuje textový formát FASTA alebo FASTQ, takže som musel najprv previesť BAM súbory do text formátu pomocou nástroja samtools. % Todo: doplniť odkaz na nástroj samtools
Po prevedení som spustil Clover, ktorý zoskupil sekvencie do jednotlivých klastrov na základe ich podobnosti.
Nastroj Clover my vratil subor ktory mal strukturu polia ktory obsahoval dvojicecelych cisel, kde prve cislo reprezentovalo ID sekvencie a druhe cislo reprezentovalo ID Sekvencie Prvej Sekvencie, do ktoreho bola sekvencia zaradena.
Zo všetkých vytvorených klastrov som vybral tie, ktoré boli dostatočne veľké na ďalšiu analýzu.
Ale Ukazalo sa ze dna v klastroch neboli zhodne dalej ako par nukleotidov, co znamenalo ze som mal dna vobec neboli klastrovane spravne.
Neskvor som zistil ze clover zarovnava iba podla zaciatku retazca a neberie do uvahy zhodu dalej v retazci, co znamenalo ze som mal v jednom klustri sekvencie ktore mali zhodu iba na zaciatku a dalej sa rozchadzali.
Jedenm klaster obsahoval viac akom 37000 sekvencií, co bolo neobvykle veľké množstvo sekvencií v jednom klustri ktore mali cez dlzku ako 43 nukleotidov zhodne, co bolo nepravdepodobné, že by sa stalo náhodou.
Vznikla teoria ze som narazil na dna Primer %Todo: doplniť odkaz na dna primer a vysvetlit ho.
Potom co som vytvoril konzesus prvich 43 nukletidov.
Na zaklade tejto informacie som naprogamoval vlastný skript v jazyku Python, ktorý pouzival kniznicu Biopython, A rozbehol som ho na vsetky vetvy a pouzil som ho na odelenie dna primerov od zvyšku sekvencií.
Potom som dane sekvencie znovu zoskupil pomocou Clover, ale tentokrát som použil iba zvyšné časti sekvencií bez dna primerov.
Data odcistene od dna primerov sa znovu zoskupili do klastrov, ktoré teraz obsahovali sekvencie, ktoré mali zhodu nielen na začiatku, ale aj ďalej v retazci.

